\chapter{多粒度感知理论与仿生轻量化方法基础}
\label{chap:theory}

\section{引言}

机载算力受限场景下的遥感视觉感知，在数学本质上是一个在严苛系统边界约束下求解结构风险最小化的带约束最优化问题（Constrained Optimization Problem, COP）。为了突破"感知精度提升必然导致算力指数级爆炸"的传统技术瓶颈，必须从底层的数学表征与跨学科的生物学先验中寻找破局点。本章将系统性地阐述遥感影像多粒度感知任务的数学空间映射模型，引入"算力-访存"双重受限的Roofline理论以严格界定边缘计算环境下的模型复杂度评价指标，并深入剖析深度视觉模型的轻量化演进机理。在此基础上，重点引入视觉仿生学与信息论（Information Theory）的理论先验，通过严密的数学推导，为后续章节中突触门控压缩、频空协同微调以及原生稀疏架构的设计提供坚实、不可辩驳的理论支撑。

\section{遥感影像多粒度视觉感知任务的数学表征与优化目标}

遥感影像的视觉感知可依据空间与特征表征的精细度划分为实例级、空间像素级与光谱像素级三个粒度层次。深度神经网络在其中的作用，本质上是去逼近一个从高维观测空间（Observation Space）到低维语义空间（Semantic Space）的高度复杂的非线性映射算子。

\subsection{实例级感知（遥感目标检测）的空间映射与损失模型}

不同于自然场景图像中目标的水平边界框（Horizontal Bounding Box, HBB），遥感影像通常由高空俯视拍摄，目标（如舰船、车辆）呈现出强烈的任意方向旋转分布。因此，实例级感知要求提取具有旋转不变性（Rotation Invariance）与等变性（Equivariance）的有向边界框（Oriented Bounding Box, OBB）。

设输入遥感影像张量为 $\mathbf{X} \in \mathbb{R}^{H \times W \times 3}$，实例级感知的目标是学习一个由参数 $\Theta_{det}$ 参数化的映射网络 $\mathcal{F}_{det}(\cdot; \Theta_{det})$，使得其输出集合逼近真实目标分布：

\begin{equation}
\hat{\mathcal{Y}} = \mathcal{F}_{det}(\mathbf{X}; \Theta_{det}) = \left\{ (\mathbf{b}_i, c_i, p_i) \right\}_{i=1}^N
\end{equation}

其中，$N$ 为影像中预测的目标总数；$c_i \in \{1, 2, \dots, C\}$ 为目标类别，具有概率 $p_i$；$\mathbf{b}_i = (x_i, y_i, w_i, h_i, \theta_i)$ 是具有五个自由度的几何描述符，分别表示第 $i$ 个预测框的中心点横纵坐标、宽度、高度及旋转角度（通常 $\theta_i \in [-\pi/2, \pi/2)$）。

该映射的优化目标（损失函数）通常解耦为分类损失 $\mathcal{L}_{cls}$ 与几何回归损失 $\mathcal{L}_{reg}$。对于遥感目标，回归损失常采用基于旋转交并比（Rotated IoU）或二维高斯分布Wasserstein距离的度量：

\begin{equation}
\mathcal{L}_{reg}(\mathbf{b}, \hat{\mathbf{b}}) = 1 - \frac{|\mathbf{b} \cap \hat{\mathbf{b}}|}{|\mathbf{b} \cup \hat{\mathbf{b}}|} \quad \text{或} \quad \mathcal{W}_2^2(\mathcal{N}(\mathbf{\mu}_1, \mathbf{\Sigma}_1), \mathcal{N}(\mathbf{\mu}_2, \mathbf{\Sigma}_2))
\end{equation}

其中，协方差矩阵 $\mathbf{\Sigma}$ 由 $(w, h, \theta)$ 通过旋转矩阵转换而来。该任务在轻量化过程中的核心理论难点是多尺度特征弥散：当网络通道数被大幅压缩时，对于高空极小目标（$w_i, h_i \to 0$），其在深层特征图上的高斯感受野响应会急剧衰减至零向量，导致梯度的梯度消失，无法进行有效的权重更新。

\subsection{空间像素级感知（语义分割）的场分布模型}

空间像素级感知旨在对遥感影像中的复杂地理拓扑结构（如受灾建筑、路网）进行密集的像素级分类，剥离背景干扰。其数学本质是一个全分辨率的空间场（Spatial Field）映射过程：

\begin{equation}
\hat{\mathbf{S}} = \mathcal{F}_{seg}(\mathbf{X}; \Theta_{seg}), \quad \hat{\mathbf{S}} \in \mathbb{R}^{H \times W \times K}
\end{equation}

对于图像空间域 $\Omega$ 中的每一个像素坐标 $(u, v)$，网络需输出一个经过Softmax归一化的概率分布向量 $\mathbf{P}(u, v) \in \Delta^{K-1}$（$K-1$维单纯形），其中 $K$ 为地物类别总数。

其理论优化目标是在经验风险最小化框架下，最小化预测分布与真实One-hot分布 $\mathbf{Y}$ 之间的交叉熵（Cross Entropy）或Kullback-Leibler散度：

\begin{equation}
\mathcal{L}_{seg} = -\frac{1}{|\Omega|} \sum_{(u,v) \in \Omega} \sum_{k=1}^K \mathbf{Y}_{k}(u,v) \log \left( \frac{\exp(\hat{\mathbf{S}}_{k}(u,v))}{\sum_{j=1}^K \exp(\hat{\mathbf{S}}_{j}(u,v))} \right)
\end{equation}

在基础视觉大模型（如Segment Anything Model, SAM）的应用范式下，该映射函数被拆解为强大的图像编码器 $\mathcal{E}$（通常为大型ViT）与轻量级的掩码解码器 $\mathcal{D}$，即 $\hat{\mathbf{S}} = \mathcal{D}(\mathcal{E}(\mathbf{X}), \text{Prompt})$。遥感影像的空间异质性使得直接利用自然图像预训练的 $\mathcal{E}$ 提取特征时，容易在拓扑边缘处产生不连续的阶跃误差（Step Error）。从信号处理的角度来看，这是由于深层网络作为低通滤波器（Low-pass Filter）滤除了遥感图像中的高频空间梯度（$\nabla \mathbf{X}$），因此亟需在模型微调中引入显式的频域高频约束。

\subsection{光谱像素级感知（高光谱分类）的高维流形模型}

高光谱图像（Hyperspectral Image, HSI）是一个包含了数百个连续狭窄波段的三维数据立方体 $\mathbf{X}_{hsi} \in \mathbb{R}^{H \times W \times B}$，其中 $B$ 为光谱波段数（通常 $B \ge 100$）。光谱像素级感知提取每一像元的光谱向量 $\mathbf{x} \in \mathbb{R}^B$，并利用空谱联合邻域 $\mathcal{N}(\mathbf{x})$ 映射至地物类别集合。

高光谱数据的核心理论挑战是统计学习理论中的**"Hughes现象"（维度灾难, Curse of Dimensionality）**。根据高维几何空间特性，当波段维度 $B$ 趋于无穷时，高维超球体的体积趋于零，样本点分布趋于超球体的表面，导致任意两点间的欧氏距离变得相近。设训练样本数为 $M$，分类器的泛化误差界（Generalization Error Bound）不仅与 $M$ 成反比，更与维度 $B$ 呈指数级正相关。如果直接在 $\mathbb{R}^B$ 空间进行全连接或3D卷积映射，分类器会迅速过拟合。

因此，从底层突破高光谱感知算力爆炸和维数灾难的关键，是设计一个数学上严密的稀疏降维变换算子 $\mathcal{T}: \mathbb{R}^B \to \mathbb{R}^d$（且 $d \ll B$），使得映射后的低维流形空间最大程度地保留原始类间的Fisher判别信息。

\section{边缘计算部署中的硬件约束与复杂度评价体系}

为了将多粒度感知算法有效部署至机载边缘控制芯片（如NVIDIA Jetson系列、FPGA），必须摒弃传统学术界单一唯精度论（仅关注mAP或OA）的评价准则，建立一套基于Roofline模型的三维硬件约束评价体系。

\subsection{时间复杂度：浮点运算量（FLOPs）}

FLOPs (Floating Point Operations) 衡量算法在执行一次前向推理时所需的理论乘加操作数。对于一个输入通道数为 $C_{in}$，输出通道数为 $C_{out}$，卷积核空间尺寸为 $K \times K$，输出特征图大小为 $H \times W$ 的标准2D卷积层，其理论计算量近似为：

\begin{equation}
\text{FLOPs}_{conv} = 2 \times K^2 \times C_{in} \times C_{out} \times H \times W
\end{equation}

对于Transformer中的多头自注意力（Multi-Head Self-Attention, MHSA）机制，假设输入序列长度为 $L = H \times W$，隐藏层维度为 $D$，其复杂度为：

\begin{equation}
\text{FLOPs}_{attention} \approx 4 L D^2 + 2 L^2 D
\end{equation}

可以看出，自注意力机制存在相对于输入分辨率 $L$ 的二次方复杂度 $\mathcal{O}(L^2)$。在机载边缘设备上，算力峰值（Peak Performance）通常被物理限制在几十至百GFLOPs量级，因此必须严格控制网络的理论算力需求。

\subsection{空间复杂度：参数规模（Params）与显存占用}

参数量（Params）衡量算法静态结构的大小。对于上述标准卷积层，其参数量为：

\begin{equation}
\text{Params}_{conv} = K^2 \times C_{in} \times C_{out} + C_{out}
\end{equation}

参数量直接决定了模型文件在机载存储器（如eMMC、Flash）中的占用大小。更为关键的是，模型在推理过程中的**激活显存（Activation Memory）**占用。高分辨率遥感影像在经过浅层网络时会产生庞大的激活张量张量池，极易触发边缘设备的显存溢出（Out of Memory, OOM）异常。

\subsection{算力-访存瓶颈：算术强度（Arithmetic Intensity）与延迟}

理论计算量低并不绝对意味着推理速度快。根据体系结构中的 Roofline模型，边缘设备的实际性能 $\mathcal{P}$（单位：GFLOPS）受到理论峰值算力 $\mathcal{P}_{max}$ 和最大内存带宽 $\beta$（单位：GB/s）的双重制约。

定义模型的算术强度 $I$ 为总浮点运算量与内存访问量（Memory Access Cost, MAC）的比值：

\begin{equation}
I = \frac{\text{FLOPs}}{\text{MAC}} \quad \text{(FLOPs/Byte)}
\end{equation}

设备的转折点算术强度为 $I_{ridge} = \mathcal{P}_{max} / \beta$。

\textbf{内存受限区（Memory-bound）：} 当 $I < I_{ridge}$ 时，模型性能 $\mathcal{P} = I \times \beta$。此时，即便模型FLOPs极低，但由于频繁读取参数和激活值耗尽了总线带宽，推理延迟（Latency）依然很高。例如，过度稀疏的非结构化剪枝网络往往落入此区域。

\textbf{计算受限区（Compute-bound）：} 当 $I \ge I_{ridge}$ 时，模型性能才能达到硬件峰值 $\mathcal{P}_{max}$。

因此，本文在进行仿生轻量化设计时，必须在降低FLOPs的同时，保证网络结构的规整性（如采用结构化剪枝或密集的原生轻量算子），以维持较高的算术强度，真正降低物理推理延迟。

\section{深度视觉模型的轻量化演进机理}

面对上述严苛的硬件边界约束，深度学习轻量化理论演化出了多种范式。本文后续的研究将基于以下三种核心机理进行深度拓展：

\subsection{模型压缩：基于泰勒展开与范数的网络稀疏化机理}

网络剪枝（Network Pruning）的核心假设是过度参数化（Over-parameterization）的深度网络中存在大量对于损失函数贡献极小的"冗余权重"。

从优化的角度来看，假设预训练模型的权重向量为 $\mathbf{W}^*$，剪枝引起权重的微小扰动 $\Delta \mathbf{W}$。将损失函数 $\mathcal{L}$ 在 $\mathbf{W}^*$ 处进行二阶泰勒展开（Taylor Expansion）：

\begin{equation}
\delta \mathcal{L} \approx \nabla \mathcal{L}^T \Delta \mathbf{W} + \frac{1}{2} \Delta \mathbf{W}^T \mathbf{H} \Delta \mathbf{W}
\end{equation}

其中，$\mathbf{H}$ 为Hessian矩阵。由于模型处于局部极小值，一阶梯度 $\nabla \mathcal{L} \approx \mathbf{0}$，因此权重去除引起的损失变化主要由Hessian项决定。由于计算全量Hessian矩阵代价过高，结构化剪枝通常采用一阶近似的启发式规则。

常见的结构化通道剪枝通过在损失函数中引入关于Batch Normalization (BN) 层缩放因子 $\gamma \in \Gamma$ 的 $\ell_1$ 范数稀疏正则化项：

\begin{equation}
\mathcal{L}_{total} = \mathcal{L}_{task}(\mathcal{F}(\mathbf{X}; \mathbf{W}, \Gamma), \mathbf{Y}) + \lambda \sum_{\gamma \in \Gamma} |\gamma|
\end{equation}

当模型收敛后，将 $|\gamma|$ 小于给定阈值的通道（及其对应的滤波器）物理移除。然而，传统的全局静态阈值剪枝存在致命缺陷：它忽略了深层网络中多尺度特征表示的流动性，粗暴的截断极其容易造成遥感小目标特征的"灾难性遗忘"。

\subsection{参数高效微调：低秩空间映射机理}

参数高效微调（Parameter-Efficient Fine-Tuning, PEFT）旨在解决拥有数十亿参数的基础大模型无法在下游任务中全量微调（Full Fine-Tuning）的困境。其底层数学机理假设：预训练大模型的过度参数化表征空间具有极低的内在秩（Intrinsic Rank）。

例如，LoRA (Low-Rank Adaptation) 机制冻结预训练权重矩阵 $\mathbf{W}_0 \in \mathbb{R}^{d \times k}$，通过引入两个极小的低秩矩阵 $\mathbf{A} \in \mathbb{R}^{d \times r}$ 和 $\mathbf{B} \in \mathbb{R}^{r \times k}$ 来近似目标域的参数更新（其中秩 $r \ll \min(d, k)$）：

\begin{equation}
\mathbf{h} = \mathbf{W}_0 \mathbf{x} + \Delta \mathbf{W} \mathbf{x} = \mathbf{W}_0 \mathbf{x} + \mathbf{A} \mathbf{B} \mathbf{x}
\end{equation}

此时，微调过程仅需优化参数集合 $\{\mathbf{A}, \mathbf{B}\}$，将梯度更新的显存开销降低了数个数量级。而Adapter机制则通过引入带有非线性激活函数 $\sigma$ 的下采样（投影矩阵 $\mathbf{W}_{down}$）和上采样（投影矩阵 $\mathbf{W}_{up}$）瓶颈结构：$\mathbf{h}' = \mathbf{h} + \mathbf{W}_{up}\sigma(\mathbf{W}_{down}\mathbf{h})$。

尽管PEFT在自然语言中极为成功，但直接将其应用于遥感图像的大模型微调时，简单的线性投影矩阵难以补偿下采样过程中丢失的空间高频梯度信息。

\subsection{轻量化算子的张量解耦机理}

利用深度可分离卷积（Depthwise Separable Convolution）等轻量化算子，其本质是将标准的致密三维卷积张量操作进行秩一分解（Rank-1 Decomposition）。它将标准卷积解耦为：

\textbf{逐通道空间卷积（Depthwise Convolution）：} 仅在二维空间域操作，参数为 $K \times K \times C_{in}$。

\textbf{逐点跨通道卷积（Pointwise Convolution）：} 即 $1 \times 1$ 卷积，仅在通道域进行线性组合。

通过这种张量解耦，计算复杂度被压缩了 $\frac{1}{C_{out}} + \frac{1}{K^2}$ 倍。然而，对于输入通道 $C_{in}$ 极大的高光谱遥感数据，这种算子级别的线性优化仍然受限于维度的绝对值，无法从根本上消除输入端的维数灾难。

\section{视觉仿生学与计算神经科学的先验基础}

自然界的生物经历了亿万年的自然选择进化，其神经中枢与视觉感知系统能够在消耗极低生物能（例如，人类大脑的峰值功率约为20W，而昆虫的大脑功耗仅在微瓦级别）的情况下，实时、极其稳健地处理充满高频噪声的外部多模态物理信号。将这种**"生物学上的节能法则与高效特征提取机理"**抽象为严格的数学模型，并跨学科地引入计算视觉与最优化求解中，是本文打破传统轻量化瓶颈的核心理论先验。

\subsection{神经突触可塑性与分形稀疏理论}

在神经生物学中，哺乳动物大脑的发育伴随着剧烈的"突触修剪（Synaptic Pruning）"。突触可塑性遵循赫布学习法则（Hebbian Learning Rule）："Fire together, wire together"。神经元的突触权重 $\mathbf{w}$ 随突触前神经元激活 $\mathbf{x}$ 和突触后神经元激活 $\mathbf{y}$ 的相关性动态演化：

\begin{equation}
\frac{d\mathbf{w}}{dt} = \eta \mathbf{y} \mathbf{x} - f(\mathbf{w}, \mathbf{y})
\end{equation}

其中 $f(\cdot)$ 为衰减项（如Oja法则中的归一化约束）。大脑利用长时程抑制（Long-Term Depression, LTD）自适应地切断那些激活频率低、携带信息少的微弱突触连接，从而在维持强大认知网络的同时，将系统熵维持在最低水平。

同时，大脑皮层的拓扑结构展现出显著的分形特征（Fractal Characteristics），即局部网络在不同的缩放尺度下表现出与整体相似的自相似性（Self-similarity），其Hausdorff维度是一个非整数。

对本文（第三章）的理论指导： 在网络剪枝中，粗暴的全局静态阈值有悖于大脑的柔性演化规律。若能模拟突触可塑性，构建具有分形特性的动态门控函数（Fractal Gating Function），在反向传播中递归地评估特征通道的重要性权重，便能指导深度模型进行自适应的稀疏化演进。这为解决大模型压缩中小目标多尺度特征极易被"剪断"的难题，提供了一种由数据驱动的柔性特征路由解法。

\subsection{人类视觉系统（HVS）的空间-频率双域协同理论}

神经生理学实验（如Hubel与Wiesel的猫视觉皮层实验）表明，灵长类动物初级视皮层（V1区）的简单细胞感受野可以用带有不同方向与频率的二维Gabor滤波器组精确拟合。更为重要的是，人类视觉系统存在对比度敏感度函数（Contrast Sensitivity Function, CSF），对空间信号呈现出强烈的带通选择性。

人类在解析复杂场景时，大脑的视觉通路潜意识地进行了信号的频空正交解耦：

\textbf{利用低频信号}（对应平缓的亮度渐变）在大脑的背侧通路（Dorsal Stream）快速构建场景的全局结构、轮廓与深度。

\textbf{利用高频信号}（对应急剧的梯度突变）在腹侧通路（Ventral Stream）敏锐捕捉物体的纹理与精细边缘。

根据二维连续傅里叶变换（2D-CFT）原理，遥感图像的空间域矩阵 $f(x,y)$ 与频域矩阵 $\mathcal{F}(u,v)$ 存在如下映射：

\begin{equation}
\mathcal{F}(u,v) = \int_{-\infty}^{\infty} \int_{-\infty}^{\infty} f(x,y) e^{-j 2\pi (ux+vy)} dx dy = R(u,v) + j I(u,v) = |\mathcal{F}(u,v)| e^{j \phi(u,v)}
\end{equation}

其中，幅值谱 $|\mathcal{F}(u,v)|$ 反映了图像的频率能量分布（低频多分布在中心，高频在边缘），而相位谱 $\phi(u,v)$ 蕴含了极其关键的空间位置与结构几何信息。

对本文（第四章）的理论指导： 视觉大模型（SAM）为了获取全局语义，其密集的下采样操作相当于一个低通滤波器，过度滤除了高频细节。借鉴HVS机制，若能在适配微调中引入快速傅里叶变换（FFT），在频域空间内显式分离出低频的几何拓扑拓扑与高频的边缘梯度，并与空域特征进行双域协同融合（Dual-domain Synergy），必将以极低的算子开销，赋予冻结权重的视觉大模型极其敏锐的破碎边缘感知能力。

\subsection{昆虫复眼光谱选择性与信息瓶颈理论}

以蜜蜂为代表的膜翅目昆虫，其大脑神经元数量极少（不足人类的百万分之一），却具备令人惊叹的光谱鉴别与视觉导航能力。昆虫神经行为学揭示，蜜蜂复眼中的光感受器并未像人工高光谱传感器那样去采集连续的数百个频段，而是仅对自然界中与其生存（如寻花蜜）息息相关的特定宽频段（紫外、蓝光、绿光）产生强烈的动作电位响应。

这种"自适应光谱选择性注意力"机制，在信息论（Information Theory）中完美契合了信息瓶颈理论（Information Bottleneck, IB）。设输入高维光谱数据为随机变量 $X$，目标标签为 $Y$，提取的紧凑特征表示为 $Z$。IB理论的目标是在最大化 $Z$ 与 $Y$ 互信息 $I(Z;Y)$ 的同时，最小化 $Z$ 与 $X$ 的互信息 $I(Z;X)$（即过滤冗余噪声）：

\begin{equation}
\min_{p(z|x)} \mathcal{L}_{IB} = I(Z;X) - \beta I(Z;Y)
\end{equation}

其中，互信息 $I(Z;Y) = \mathcal{H}(Y) - \mathcal{H}(Y|Z)$，$\mathcal{H}(\cdot)$ 为香农熵（Shannon Entropy）。蜜蜂的视觉神经系统正是自然界中实现该优化目标的完美物理实体：它通过硬连线的神经通路，直接在传感器端掐断了低信息熵的频段，仅保留判别性极高的"黄金波段"。

对本文（第五章）的理论指导： 针对高光谱遥感数据指数级爆炸的计算维数灾难，必须摒弃深度学习惯用的"先全量编码、再在网络后端降维"的重型范式。通过将其抽象为极简数学模型，在神经网络的最前端植入基于互信息或交叉熵梯度的"动态门控选择器"，强行模拟蜜蜂的光谱稀疏截断机制，是实现"计算前置"，将前向推理算力下限压缩至 KB 级别（$<100$ KFLOPs）的最根本理论途径。

\section{本章小结}

本章从数学空间映射、硬件物理约束以及生物学底层法则的多维交叉视角，为整篇学位论文的算法设计构筑了严密、完备的理论地基。首先，通过构建带约束的优化目标，严格定义了实例级检测、空间像素级分割以及光谱像素级分类三种多粒度任务的函数模型；引入Roofline评价体系，揭示了机载边缘设备中FLOPs、参数量与访存带宽限制算力发挥的真实根源。其次，利用泰勒展开、矩阵低秩近似以及张量解耦等数学工具，深度剖析了现行模型压缩、微调及算子轻量化理论的运行机理及其在面临极端算力与遥感复杂特征冲突时的技术天花板。最后，创造性地梳理了神经突触可塑性与分形理论、人类视觉频空双域协同理论以及昆虫复眼信息瓶颈先验。这些跨越计算神经科学与信息论的仿生法则，不仅在哲学层面揭示了低功耗、高鲁棒智能感知系统的终极形态，更在随后的数学演绎中，为突破遥感轻量化解译的技术死局提供了直接的方法论钥匙。

